% ===========================================================================
%  Chapter 8: Discussion
%  Dual-Target RL Molecular Generation — Workshop Report
% ===========================================================================

% >>> PLACEHOLDER: this section is populated after pipeline completion. <<<
% Subsection structure below defines the narrative arc; each subsection
% should be expanded with quantitative references to the Results tables.

\subsection{Interpretation of Primary Results}

% >>> PLACEHOLDER: interpret hr_both differences, effect sizes, CI. <<<
% Relate N/5 success rate to pre-registered criteria (Section 4).
% Discuss whether the effect is consistent across homo vs hetero pairs.
% >>> END PLACEHOLDER <<<

\subsection{Comparison with Prior Work}

No prior work has systematically explored the continuous weight space
from selective ($w_A = 1, w_B = -1$) through single-target
($w_A = 1, w_B = 0$) to dual-target ($w_A = 1, w_B = 1$) for
structure-based docking as the RL reward for two protein
targets~\citep{fromer2023moo, loeffler2024reinvent4}.
Existing multi-objective RL generators either optimize molecular
\emph{properties} (QED, SA, logP) rather than binding affinity to
two specific proteins~\citep{jain2023mogfn, wang2024scarlp},
or use predictor-based rewards instead of
docking~\citep{ren2024polygon, zhang2024clm}.

The closest comparisons are:
\begin{itemize}
\item \textbf{PMMG}~\citep{pmmg2025}: Pareto MCTS for dual EGFR/HER2
  docking, but without weight-space exploration or orthogonal evaluation.
\item \textbf{AIxFuse}~\citep{chen2024aixfuse}: pharmacophore-based
  dual-target design for GSK3$\beta$/JNK3, but no selectivity spectrum.
\item \textbf{Isigkeit et al.}~\citep{isigkeit2024clm}: chemical language
  model fine-tuning for dual ligands with experimental validation, but
  ligand-based (no structural docking reward).
\end{itemize}

% >>> PLACEHOLDER: compare quantitative results to PMMG/POLYGON success rates. <<<
% >>> END PLACEHOLDER <<<

\subsection{Orthogonal Evaluation: UniDock vs.\ Boltz-2}

A key design choice is the decoupling of the training objective (UniDock
docking scores $d_A, d_B$) from the evaluation metric (Boltz-2 affinity
$a_A, a_B$)~\citep{passaro2025boltz2}.  Most generative studies evaluate
molecules with the same scorer used as the RL reward, creating circular
evaluation~\citep{abdelsayed2025polypharmacology}.
Our design mitigates this by using two fundamentally different
methods---physics-based empirical scoring versus a transformer-based
learned surrogate.

% >>> PLACEHOLDER: report per-target Spearman rho and discuss concordance. <<<
% If rho ~ 0 for any target, flag hit-rate conclusions for that pair.
% >>> END PLACEHOLDER <<<

\citet{bret2026boltz2docking} recently showed that Boltz-2 affinity
predictions may be insensitive to binding-site mutations, reinforcing the
need for multiple independent evaluation layers.  Where docking and Boltz-2
agree on the direction of the \texttt{double} advantage, confidence in the
result is substantially higher.

\subsection{Weight-Space Spectrum}

The six reward modes map the trade-off between target selectivity and
polypharmacology onto a single parametric axis.  This is, to our knowledge,
the first demonstration that these two traditionally separate research
fields correspond to different points on a continuous reward-weight
scale within the same RL architecture.

% >>> PLACEHOLDER: summarize trends across the 6 modes. <<<
% Does selective_a mode actually produce A-selective molecules?
% Does linear_equal differ from worst-case?
% >>> END PLACEHOLDER <<<

\subsection{Limitations in Context}

The limitations enumerated in Section~\ref{sec:limitations} should be
weighed against three mitigating factors:
(1)~all comparisons are \emph{relative} (double vs.\ single under
identical conditions), so absolute scorer calibration is less critical;
(2)~pre-registration of pairs, thresholds, and success criteria
prevents post-hoc cherry-picking;
(3)~the conjunctive significance criterion (FDR $+$ effect size $+$
CI exclusion) is deliberately conservative.

The $m = 5$ seed limitation remains the principal statistical weakness.
Future work should increase the seed count to $\geq 10$ and consider
wild cluster bootstrap~\citep{cameron2008bootstrap} for improved
small-cluster inference.
